%
% 3_Resultante.tex -- Abschnitt 3 Resultante 
%
% (c) 2020 Prof Dr Andreas Müller, Hochschule Rapperswil
%
% !TEX root = ../../buch.tex
% !TEX encoding = UTF-8
%
\section{Resultante
\label{resultante:section:resultante}}
\kopfrechts{Resultante}

Die Resultante ist ein Werkzeug der kommutativen Algebra, um das Vorhanden sein gemeinsamer Nullstellen zweier Polynome zu prüfen.
Die Resultante kann dazu verwendet werden, nacheinander die Variablen eines Gleichungssystems mit mehreren Variablen zu eliminieren.
Zu diesem Zweck wurde die Resultante im \(19.\) Jahrhundert untersucht.
In modernen CAS werden Resultanten unteranderem benutzt um aus einer vorher bestimmten Gröbner-Basis auf die Lösung eines Gleichungssystems zu schliessen. \cite{resultante:resultante}

\subsection*{Nützliche Eigenschaften
\label{resultante:subsection:eigenschaften}}

Die erste nützliche Eigenschaft der Resultante zweier Polynome besteht darin, das sie sich in Bezug auf deren Nullstellen ausdrücken lässt.
Die Resultante bietet einen Weg zur parmetrischen Berechnung der Nullstellen.
